\begin{abstract}
    % Large vision models are increasingly strong at vision tasks in natural images, however they still struggle in low signal-to-noise ratio (SNR) videos, such as underwater sonar, ultrasound, and microscopy videos.
    Low signal-to-noise ratio (SNR) videos—such as those from underwater sonar, ultrasound, and microscopy—pose significant challenges for computer vision models, particularly when paired clean imagery for denoising is unavailable.
    We present Spatiotemporal Augmentations and denoising in Video for Downstream Tasks (\NAME{}), a novel self-supervised method that denoises low-SNR sensor videos using only raw noisy data. By leveraging distinctions between foreground and background motion and exaggerating objects with stronger motion signal, \NAME{} enhances foreground object visibility and reduces background and camera noise without requiring clean video. \NAME{} has a set of architectural optimizations that lead to faster throughput, training, and inference than existing deep learning methods. We also introduce a new denoising metric, FBD, which indicates foreground-background divergence for detection datasets without requiring clean imagery. Our approach achieves state-of-the-art results for classification, detection, tracking, and counting tasks and it does so with fewer training resource requirements than existing deep-learning-based denoising methods. Project page \href{https://suzanne-stathatos.github.io/SAVeD}{here}, Code: \href{https://github.com/suzanne-stathatos/SAVeD}{https://github.com/suzanne-stathatos/SAVeD}.
\end{abstract}

\iffalse
\begin{abstract}
    Foundation models excel at vision tasks in natural images but fail in low signal-to-noise ratio (SNR) videos, such as underwater sonar, ultrasound, and microscopy. We introduce Spatiotemporal Augmentations and denoising in Video for Downstream Tasks (\NAME{}), a self-supervised method that denoises low-SNR sensor videos and is trained using only the raw noisy data. By leveraging differences in foreground and background motion, \NAME{} enhances object visibility using an encoder-decoder with a temporal bottleneck. \NAME{} also has a set of architectural optimizations that lead to much faster training than existing deep learning methods. We also introduce a new denoising metric, FBD, which requires no clean ground-truth imagery. Our approach achieves SOTA results for classification, detection, tracking, and counting performance and it does so with fewer resource requirements than existing deep-learning-based denoising methods.
\end{abstract}

\begin{abstract}
With the rise of foundational models in vision, tasks like classification, detection, and object tracking in natural images increasingly work out of the box.
These foundation models do not work for sensory videos with low signal-to-noise ratios, such as underwater sonar of swimming fish, ultrasonic scans of breast tissue, or microscopy videos of moving cells.
We present Spatiotemporal Video Denoising for Downstream Tasks (\NAME{}), a novel method that uses self-supervised learning to denoise sensor videos without requiring noise-free video.
We show how \NAME{} improves performance on downstream tasks such as classification, detection, tracking, and counting.
Sensor data, such as underwater sonar or medical ultrasound, presents a few additional challenges from above-ground RGB video. 
First, there's not as much training data available as for RGB. 
%Second, in underwater sonar videos, the background -- including water flow, shadowing, and environmental factors -- does not stay stationary, objects of interest range in scale within the frame, making small objects difficult to identify without their motion. 
% Second, in underwater sonar videos, the dynamic background and varying object scales make small objects hard to identify without motion.
% In a stationary frame, the moving objects do not have identifiable features, but we observe that foreground objects have a motion signature that is very different from the background motion. In ultrasound videos, the camera is also capturing a liquid environment, however, there is camera motion. In these videos, we aim to enhance the characteristics of the lesions and diseases and denoising the speckle from the video.
In underwater sonar and ultrasound videos, dynamic backgrounds, scale variations, and camera motion obscure small objects and features. However, foreground motion differs from background motion, which we leverage to denoise.
 \NAME{} uses an encoder-decoder architecture with a temporal bottleneck to reconstruct the spatiotemporal difference between video frames, capturing motion signatures and exaggerating objects of interest. 
 By focusing on regions of movement and filtering out movement with minimal signal, our approach works directly on noisy video without requiring any form of annotation. 
% Our approach works directly on noisy video without requiring any form of annotation. 
We show that denoising the video frames improves downstream object detection, tracking, counting, and classification performance. 
Finally, we show that our method beats other state-of-the-art denoising methods while requiring less resources.

\end{abstract}
\fi


% For the denoiser, put noisy vs. denoised to see qualitatively what's happening -- in 2 rows, explicitly separate the top (setup) and bottom (applications) and put small titles over each \ie detecting fish, categorizing lung disease, categorizing breast cancer

% For all the fish ones, make them match \ie frame from denoiser and frame to downstream 

    